O escopo deste trabalho revelou-se seu maior desafio. Embora a base do SNIRH seja ampla, apresentou diversos problemas
de qualidade, inconsistências e falhas nas séries temporais. O esforço necessário para corrigir e padronizar esses dados
consumiu uma parcela significativa do tempo, que poderia ter sido direcionada para etapas centrais da modelagem, como
engenharia de atributos mais profunda, testes adicionais de arquiteturas, refinamento sistemático das predições, ou a
inclusão de outras fontes de dados.

A adoção de uma arquitetura em paralelo também se mostrou complexa. Essa escolha foi motivada por um cenário realista,
no qual um sistema operacional de previsão de cheias não pode depender de um único modelo, mas deve integrar diferentes
arquiteturas e comparar múltiplos cenários, potencialmente por meio de um \textit{ensemble}. Contudo, treinar
simultaneamente modelos distintos com entradas idênticas e configurações otimizadas, excedeu as capacidades de uma
máquina comercial. Um ambiente computacional mais robusto certamente produziria resultados superiores. Mais memória RAM
permitiria ampliar o conjunto de atributos selecionados; uma GPU de maior capacidade proporcionaria paralelismo mais
efetivo e tempos de treinamento substancialmente menores. Ainda assim, por tratar-se de um estudo de caso, foi
necessário trabalhar com os recursos disponíveis.

Os resultados obtidos foram satisfatórios, já que todos os modelos superaram a baseline e apresentaram capacidade
preditiva consistente. A reprodução dos padrões gerais de variação do nível d'água também foi robusta. Embora nenhum
modelo tenha capturado adequadamente os picos extremos das enchentes, isso era esperado, visto que tal evento foi
extraordinário e não possuia análogos nas séries históricas utilizadas para o treinamento. Ainda assim, a habilidade dos
modelos em identificar aumentos e quedas repentinas mostrou-se notável.

Entre as abordagens avaliadas, os modelos baseados em \textit{Gradient Boosting} demonstraram maior potencial,
especialmente pela combinação de desempenho e velocidade de treinamento. Além disso, esses modelos oferecem margem para
aperfeiçoamentos adicionais, notadamente no tratamento do escalonamento das variáveis, cuja aplicação pode interagir de
forma indesejada com métodos baseados em árvores, uma vez que seus mecanismos internos podem ser capazes de lidar com
este problema.